% Generated draft: requires mathematical and source review.
\documentclass[11pt]{article}
\usepackage[margin=25mm]{geometry}
\usepackage{fontspec}
\setmainfont{Latin Modern Roman}
\usepackage{amsmath,amssymb,mathtools,unicode-math}
\setmathfont{Latin Modern Math}
\usepackage{enumitem,needspace,xurl,hyperref}
\setlength{\parindent}{0pt}
\setlength{\parskip}{5pt}
\setlength{\emergencystretch}{3em}
\newcommand{\N}{\mathbb{N}}
\newcommand{\Z}{\mathbb{Z}}
\newcommand{\Q}{\mathbb{Q}}
\newcommand{\R}{\mathbb{R}}
\newcommand{\C}{\mathbb{C}}
\newcommand{\F}{\mathbb{F}}
\newcommand{\A}{\mathbb{A}}
\newcommand{\PP}{\mathbb{P}}
\newcommand{\E}{\mathbb{E}}
\newcommand{\eps}{\varepsilon}
\newcommand{\dd}{\,\mathrm{d}}
\begin{document}
\section*{UnsolvedMath ID: 30001403}
\textbf{Source problem code:} catalogue entry
\subsection{Definitions and mathematical statement}
Let \(K\) be a field, regarded as a first-order structure in the language intended by the source. No restriction on the characteristic is stated. Its complete theory \(T=\operatorname{Th}(K)\) is called stable if no formula \(\varphi(x;y)\) has the order property: there are no sequences \((a_i)_{i<\omega}\) and \((b_j)_{j<\omega}\) in a model of \(T\) such that \(\varphi(a_i;b_j)\) holds exactly when \(i<j\).

For context, the standard weight of a complete type \(p\) in a stable theory is the supremum of the cardinalities \(|I|\) for which there exist a nonforking extension of \(p\) over a parameter set \(B\), a realization \(a\) of that extension, and a family \((b_i)_{i\in I}\) independent over \(B\), such that \(\operatorname{tp}(a/Bb_i)\) forks over \(B\) for every \(i\in I\). A type has finite weight when this supremum is finite. This definition alone does not specify what the source means by a field having finite weight: the record does not identify the relevant type or collection of types. The record also does not define \emph{strongly stable}. These qualifications therefore remain unresolved; they are retained as source predicates in the statements below.

A field \(K\) is \emph{separably closed} if it has no proper finite separable algebraic extension. It is \emph{algebraically closed} if every nonconstant polynomial in \(K[X]\) has a root in \(K\).

\textbf{Conjecture 1, as supplied.} For every infinite field \(K\),
\[
\bigl(K\text{ is stable and has finite weight in the source's intended sense}\bigr)
\ \Longrightarrow\ K\text{ is separably closed}.
\]
\textbf{Conjecture 2, as supplied.} For every infinite field \(K\),
\[
K\text{ is strongly stable in the source's intended sense}
\ \Longrightarrow\ K\text{ is algebraically closed}.
\]
The intended language and the meanings of the two weight-related qualifications require confirmation from the cited sources before these become fully specified first-order model-theoretic conjectures. The displayed implications preserve the supplied targets without resolving those omissions by assumption.
\subsection{Short English statement}
The source proposes two closure conjectures: every infinite stable field of finite weight is separably closed, and every infinite strongly stable field is algebraically closed. The intended language, the field-level finite-weight convention, and the meaning of strongly stable require confirmation. The supplied record labels the problem open; this draft does not independently verify that status.
\subsection{Sources}
\begin{itemize}
\item Local UnsolvedMath registry, permanent ID \(30001403\), source problem code \texttt{catalogue entry}, in the Oberwolfach Reports Open Problems collection. The supplied original and clean statements coincide; the registry marks the statement as exact and retained byte-for-byte. No individual proposer is named; the supplied proposed year is 2010.
\item \emph{Model Theory: Around Valued Fields and Dependent Theories} (2010), Oberwolfach Reports, pp.\ 5--53, DOI \texttt{10.4171/owr/2010/01}. Supplied source report and literature-source link labelled ``original Oberwolfach problem report'': \url{https://doi.org/10.4171/owr/2010/01}. The supplied background describes Oberwolfach Reports as documenting mathematical workshops and their open-problem discussions.
\item Krupiński and Pillay, supplied literature-source attribution: \url{https://doi.org/10.1017/S1474748010000228}. No paper title or more specific bibliographic details are supplied.
\item Supplied literature assessment, checked on \texttt{2026-08-21}: the cited paper is reported to prove the weight-one case and to record the finite-weight assertion as Conjecture 0.3; the assessment further reports that later work still treats the general stable-fields assertion as a conjecture. The supplied status is ``open''. This is supplied dated curation derived from the OWR record, not an independently verified literature assessment. The record expressly requests independent verification of time-sensitive status and cited claims.
\end{itemize}
\end{document}
